% TEMPLATE-MILC-v3.tex - plantilla maestra UNICA de las guias MILC v3.
%
% La usan las tres familias de guias, cada una con su driver:
%   · Grados 6-11  -> scripts/build-guias-g11.py
%   · Bebras       -> scripts/build-guias-bebras.py
%   · Semillero    -> scripts/build-guias-semillero.py
%
% Antes habia tres copias casi identicas (~400 lineas iguales, 6 distintas):
% cada mejora habia que aplicarla tres veces, y en la practica no se hacia
% (el motor de recursos vivia solo en dos de ellas). Lo que varia entre
% familias esta parametrizado en 6 placeholders que cada driver llena
% (se nombran aqui SIN su sintaxis real: el builder reemplaza por texto plano,
% tambien dentro de los comentarios, y un valor multilinea rompe el preambulo):
%
%   HEADER_IZQ        encabezado izquierdo de pagina
%   HEADER_DER        encabezado derecho de pagina
%   PORTADA_ETIQUETA  rotulo sobre el numero grande (GRADO/MOMENTO/MODULO)
%   PORTADA_SERIE     linea de serie de la portada
%   PORTADA_ID        identificador de la guia en la portada
%   PORTADA_CIRCULO   el circulito de grado (°), o vacio si no aplica
%
% El resto de claves las documenta content/guias/_SCHEMA.md. El script valida
% que no queden placeholders sin reemplazar antes de escribir el archivo final.

\documentclass[11pt,letterpaper]{article}
\usepackage[margin=2.0cm]{geometry}
\usepackage[table]{xcolor}
\usepackage{fontspec}
\usepackage[spanish]{babel}
\usepackage{tikz}
% Librerías TikZ para los diagramas de las guías (bloque `recursos:`):
%  · babel  → babel en español vuelve activos caracteres como «>», y TikZ los usa
%             en sus opciones (>=stealth, ->). Sin esta librería, cualquier
%             diagrama con flechas revienta con «\language@active@arg> has an
%             extra }». Debe ir ANTES de que se dibuje nada.
%  · positioning        → sintaxis «below left=2pt and 3pt».
%  · decorations.pathreplacing → llaves ({brace}) para señalar rangos.
\usetikzlibrary{babel,positioning,decorations.pathreplacing}
\usepackage{graphicx}
% Circuitos: el trabajo de electrónica se hace hoy con micro:bit y sus sensores
% integrados (MakeCode, sin cableado), así que no se necesitan esquemáticos.
% Si algún día entran montajes con protoboard, basta descomentar esta línea:
% los diagramas .tex/.tikz declarados en `recursos:` ya se incrustan solos.
% \usepackage{circuitikz}
\usepackage[most]{tcolorbox}
\usepackage{tabularx}
\usepackage{array}
\usepackage{fancyhdr}
\usepackage{titlesec}
\usepackage[document]{ragged2e}  % bandera a la derecha en todo el cuerpo
\usepackage{enumitem}
\usepackage{microtype}

% Helvetica Neue no trae la flecha U+2192, asi que cada «→» escrito literalmente
% en el YAML se compilaba a NADA, en silencio (462 flechas en 61 guias: rutas del
% tipo «paleta → tipografia → lockup» salian rotas). Mapeamos el caracter a su
% version matematica, que si tiene glifo. Asi el autor puede seguir escribiendo
% «→» comodamente en el YAML.
\usepackage{newunicodechar}
\newunicodechar{→}{\ensuremath{\rightarrow}}
% Mismo problema con los simbolos de marca y vinetas: el «✓» de «una palomita ✓»
% salia como cuadrito vacio en el PDF de todas las guias que lo usaban.
\usepackage{pifont}
\newunicodechar{✓}{\ding{51}}
\newunicodechar{✗}{\ding{55}}
\newunicodechar{✔}{\ding{52}}
\newunicodechar{•}{\textbullet}
\newunicodechar{≥}{\ensuremath{\geq}}
\newunicodechar{≤}{\ensuremath{\leq}}

\setmainfont{Helvetica Neue}
\setsansfont{Helvetica Neue}
\setlength{\parindent}{0pt}
\setlength{\parskip}{6pt}
% Sin particion de palabras: en 6.º-7.º una palabra cortada por guion al final
% de linea («Comuni-cacion») frena la lectura mucho mas que una linea corta.
\hyphenpenalty=10000
\exhyphenpenalty=10000
\pretolerance=2000
\tolerance=3000
\setlength{\emergencystretch}{3em}
\linespread{1.10}
\renewcommand{\arraystretch}{1.25}
\setlist{leftmargin=5mm,itemsep=1mm,topsep=1mm}
\newcolumntype{Y}{>{\RaggedRight\arraybackslash}X}

\definecolor{milcMagenta}{HTML}{E6007E}
\definecolor{milcVino}{HTML}{5A0038}
\definecolor{milcTurquesa}{HTML}{0093A5}
\definecolor{milcVerde}{HTML}{58A923}
\definecolor{milcMostaza}{HTML}{E5B400}
\definecolor{milcOcre}{HTML}{B86600}
\definecolor{milcNegro}{HTML}{241816}
\definecolor{milcGris}{HTML}{F7F7F4}
\definecolor{milcLinea}{HTML}{E8E2DD}
\definecolor{milcGradePrimary}{HTML}{84CC16}
\definecolor{milcGradeDark}{HTML}{4D7C0F}
\definecolor{milcGradeSoft}{HTML}{ECFCCB}
\definecolor{milcGradeAccent}{HTML}{CFE7FF}
\definecolor{milcGradeText}{HTML}{FFFFFF}
\color{milcNegro}

\pagestyle{fancy}
\fancyhf{}
\lhead{\small Tecnología e Informática · Grado 7}
\rhead{\small Guía 30 · MILC v3}
\cfoot{\small\thepage}
\renewcommand{\headrulewidth}{0pt}

\newtcolorbox{titlebox}[1]{
  enhanced,colback=#1,colframe=#1,arc=7mm,boxrule=0pt,
  left=7mm,right=7mm,top=3mm,bottom=3mm,
  before skip=8pt,after skip=8pt,
  fontupper=\sffamily\bfseries\Large\color{white}
}

\newtcolorbox{softbox}[3]{
  enhanced,colback=#2,colframe=#1,borderline west={4pt}{0pt}{#1},
  arc=2mm,boxrule=.4pt,left=4mm,right=4mm,top=3mm,bottom=3mm,
  before skip=6pt,after skip=8pt,title={#3},coltitle=white,
  colbacktitle=#1,fonttitle=\sffamily\bfseries,fontupper=\RaggedRight,
  attach boxed title to top left={xshift=3mm,yshift=-2mm},
  boxed title style={arc=2mm, boxrule=0pt}
}

\newtcolorbox{drawbox}[2]{
  enhanced,colback=white,colframe=#1,arc=4mm,boxrule=1pt,
  left=5mm,right=5mm,top=4mm,bottom=4mm,before skip=8pt,after skip=8pt,
  title={#2},coltitle=white,colbacktitle=#1,fonttitle=\sffamily\bfseries,
  fontupper=\RaggedRight,
  attach boxed title to top left={xshift=4mm,yshift=-2mm},
  boxed title style={arc=3mm, boxrule=0pt}
}

\newtcolorbox{infoband}[1]{
  enhanced,colback=#1!8,colframe=#1!50,arc=2mm,boxrule=.5pt,
  left=4mm,right=4mm,top=2mm,bottom=2mm,before skip=4pt,after skip=4pt,
  fontupper=\small\RaggedRight
}

% Puente narrativo entre secciones (contrato editorial Conéctate).
% Da continuidad: "Acabas de X. Ahora vas a Y porque Z."
% Visualmente: banda izquierda en color del grado, texto en itálica pequeña.
\newtcolorbox{puentebox}{
  enhanced,colback=milcGradeSoft,colframe=milcGradeDark,
  borderline west={3pt}{0pt}{milcGradeDark},
  arc=0mm,boxrule=0pt,left=5mm,right=4mm,top=2.5mm,bottom=2.5mm,
  before skip=4pt,after skip=8pt,
  fontupper=\itshape\small\color{milcNegro}\RaggedRight
}
% La flecha va en modo matematico a proposito: Helvetica Neue NO tiene el glifo
% U+2192, asi que \textrightarrow se compilaba a NADA («Missing character: There
% is no → in font Helvetica Neue») y el puente salia sin flecha. En modo
% matematico la flecha viene de la fuente matematica y si se dibuja.
\newcommand{\puente}[1]{%
  \begin{puentebox}%
    \textcolor{milcGradeDark}{$\rightarrow$}\hspace{2mm}#1%
  \end{puentebox}%
}

\newcommand{\linead}[1]{\noindent\color{milcLinea}\rule{#1}{0.5pt}\color{milcNegro}}
\newcommand{\respuesta}[1]{\par\vspace{2mm}\foreach \n in {1,...,#1}{\noindent\color{milcLinea}\rule{\linewidth}{0.5pt}\par\vspace{5mm}}\color{milcNegro}}
\newcommand{\checkbox}{\textcolor{milcGradeDark}{\fbox{\phantom{x}}}\hspace{5pt}}

\newcommand{\phasecard}[4]{
\begin{tcolorbox}[enhanced,colback=#3,colframe=#2,arc=3mm,boxrule=.5pt,height=3.05cm,valign=center,halign=center,left=1.5mm,right=1.5mm,top=2mm,bottom=2mm]
{\sffamily\bfseries\fontsize{22}{24}\selectfont\textcolor{#2}{#1}}\par
{\sffamily\bfseries\small #4}
\end{tcolorbox}
}

% ── Piezas del rediseno 2026 (legibilidad 6.º-11.º) ───────────────────
% Banda de estacion: reemplaza el titlebox macizo. Titulo a la izquierda,
% dato practico (tiempo/modalidad) a la derecha, en una sola linea.
\newcommand{\milcestacion}[2]{%
\begin{tcolorbox}[enhanced,colback=#1,colframe=#1,arc=2mm,boxrule=0pt,
  left=5mm,right=5mm,top=2.6mm,bottom=2.6mm,before skip=11pt,after skip=7pt]
{\sffamily\bfseries\fontsize{15}{18}\selectfont\color{white}#2}
\end{tcolorbox}}

% Tarjeta de paso numerada: sustituye las tablas «Pilar / Aplicacion», que
% eran formato de consulta docente, no de estudiante. El numero va en un
% circulo sobre el borde para que la secuencia se lea de un vistazo.
\newcommand{\milcpaso}[3]{%
\begin{tcolorbox}[enhanced,colback=white,colframe=milcLinea,arc=1.5mm,boxrule=.9pt,
  left=14mm,right=4mm,top=3mm,bottom=3mm,before skip=4pt,after skip=4pt,
  fontupper=\RaggedRight,
  overlay={\node[anchor=north west,circle,fill=#2,text=white,inner sep=0pt,
    minimum size=7.4mm,font=\sffamily\bfseries\small]
    at ([xshift=3.6mm,yshift=-3.2mm]frame.north west) {#1};}]
#3
\end{tcolorbox}}

% Casilla marcable de tamano util con lapiz (la anterior era \fbox{\phantom{x}},
% de alto variable segun el texto que la seguia).
\newcommand{\milcmarca}{\framebox[3.8mm]{\rule{0pt}{3.4mm}}}

% Fila marcable de lista suelta (checks de la estacion 1).
\newcommand{\milccheckrow}[1]{%
\par\vspace{1.8mm}\noindent
\makebox[6.4mm][l]{\raisebox{-.55mm}{\textcolor{milcNegro}{\milcmarca}}}%
\parbox[t]{\dimexpr\linewidth-6.4mm\relax}{\RaggedRight #1}\par}

% Celda marcable dentro de la rubrica (tabla de autoevaluacion). Misma
% sangria francesa, pero pensada para vivir dentro de una columna X.
\newcommand{\milcopcion}[1]{%
\makebox[5.2mm][l]{\raisebox{-.55mm}{\textcolor{milcNegro}{\milcmarca}}}%
\parbox[t]{\dimexpr\linewidth-5.2mm\relax}{\RaggedRight #1}}

% ── Recursos visuales de la guía ──────────────────────────────────────
% Declarados en el bloque `recursos:` del YAML y emitidos por el builder.
% \guiaFigura   → imagen rasterizada o PDF (foto, carta celeste, montaje).
% \guiaDiagrama → fuente TikZ/circuitikz (.tex) incrustada como vector:
%                 es la vía para circuitos y esquemáticos editables.
% #1 = ruta relativa al .tex · #2 = pie (puede ir vacío)
\newcommand{\guiaFigura}[2]{%
\begin{center}
\includegraphics[width=0.88\linewidth,height=0.38\textheight,keepaspectratio]{#1}\\[1.5mm]
{\footnotesize\itshape\textcolor{milcNegro!75}{#2}}
\end{center}
}

\newcommand{\guiaDiagrama}[2]{%
\begin{center}
\input{#1}\\[1.5mm]
{\footnotesize\itshape\textcolor{milcNegro!75}{#2}}
\end{center}
}

\begin{document}
\thispagestyle{empty}

% ── Portada ───────────────────────────────────────────────────────────
% Antes: un rectangulo de color a pagina completa. Se veia bien en pantalla,
% pero en la fotocopiadora del colegio se comia el toner y salia gris sucio.
% Ahora la banda ocupa el tercio superior y el resto va en blanco.
\begin{tikzpicture}[remember picture,overlay]
  \path[fill=milcGradePrimary]
    (current page.north west) rectangle ([yshift=-10.6cm]current page.north east);

  \node[anchor=north west,text=milcGradeText] at ([xshift=2.0cm,yshift=-1.55cm]current page.north west)
    {\sffamily\bfseries\fontsize{12}{14}\selectfont GRADO};
  \node[anchor=north west,text=milcGradeText,opacity=.85] at ([xshift=2.0cm,yshift=-2.35cm]current page.north west)
    {\sffamily\bfseries\fontsize{8.5}{10}\selectfont SERIE GUÍAS · TECNOLOGÍA E INFORMÁTICA};
  \node[anchor=north east,text=milcGradeText,opacity=.85,align=right] at ([xshift=-2.0cm,yshift=-1.55cm]current page.north east)
    {\sffamily\bfseries\fontsize{9}{12}\selectfont I.E. Sor María Juliana\\Cartago, Valle del Cauca};

  \node[anchor=west,text=milcGradeText] (gradonum) at ([xshift=1.85cm,yshift=-7.1cm]current page.north west)
    {\sffamily\bfseries\fontsize{112}{112}\selectfont 7};
  % El circulito de grado (°) solo aplica a las guias de grado («11°»); en
  % Bebras y Semillero el numero grande es un momento/modulo, no un grado.
  % Se posiciona relativo al nodo `gradonum`, no en coordenadas absolutas.
  \draw[milcGradeText,line width=1.6pt]
    ([xshift=2.0mm,yshift=-4.5mm]gradonum.north east) circle (.15cm);

  \node[anchor=west,text=milcGradeText,text width=12.3cm,align=left] at ([xshift=6.5cm,yshift=-6.6cm]current page.north west)
    {
     {\sffamily\bfseries\fontsize{9}{11}\selectfont\MakeUppercase{Período 3 · Inteligencia Artificial}}\\[3mm]
     {\sffamily\bfseries\fontsize{21}{25}\selectfont Cosecha\\del periodo 3 ---\\mi proyecto\\con IA\\responsable}\\[3.5mm]
     {\sffamily\bfseries\fontsize{15}{18}\selectfont Guía 30-7-TIC}
    };
\end{tikzpicture}

\vspace*{9.4cm}
\vfill

{\sffamily\bfseries\fontsize{8}{10}\selectfont\textcolor{milcGradeDark}{LO QUE TE LLEVAS HOY}}

\begin{tcolorbox}[enhanced,colback=white,colframe=milcNegro,arc=2mm,boxrule=1.6pt,
  left=5mm,right=5mm,top=4mm,bottom=4mm,before skip=3pt,after skip=12pt,
  fontupper=\RaggedRight\fontsize{11}{15.5}\selectfont]
Un \textbf{proyecto completo} que use IA con criterio: \textbf{(a) tema escogido} (informativo, no opinión); \textbf{(b) documento final} (Word o presentación) con tu voz dominante, IA como asistente declarado; \textbf{(c) anexo de transparencia}: 3-5 prompts CTRF/avanzados usados + capturas de las respuestas + reflexión sobre cómo verificaste/adaptaste; \textbf{(d) sustentación oral} de 5 minutos al profe; \textbf{(e) autoevaluación} con rúbrica de 10 criterios y \textbf{cierre del año}: ¿qué tipo de estudiante de IA soy ahora vs hace 10 sesiones?.
\end{tcolorbox}

\begin{tcolorbox}[enhanced,colback=milcGris,colframe=milcGris,arc=2mm,boxrule=0pt,
  left=5mm,right=5mm,top=3.5mm,bottom=3.5mm,after skip=12pt]
{\sffamily\bfseries\fontsize{8}{10}\selectfont\textcolor{milcNegro!55}{ESTA GUÍA ES DE}}\\[3mm]
\begin{tabularx}{\linewidth}{X p{3.2cm} p{3.2cm}}
\linead{\linewidth} & \linead{3.0cm} & \linead{3.0cm}\\[-1mm]
{\fontsize{8.5}{10}\selectfont\textcolor{milcNegro!55}{Nombre completo}} &
{\fontsize{8.5}{10}\selectfont\textcolor{milcNegro!55}{Grupo}} &
{\fontsize{8.5}{10}\selectfont\textcolor{milcNegro!55}{Fecha}}\\
\end{tabularx}
\end{tcolorbox}

\vfill

\noindent\textcolor{milcLinea}{\rule{\linewidth}{1pt}}\\[2mm]
{\sffamily\bfseries\fontsize{9.5}{12}\selectfont Autor: Dr. Álvaro Cárdenas Orozco}\\[1mm]
{\sffamily\fontsize{8.5}{11}\selectfont\textcolor{milcNegro!55}{Metodología MILC · apertura ancestral · pensamiento computacional · triángulo Dussel–estoicismo–Floridi}}

\newpage

\section*{Apertura: Cosecha del periodo 3 --- mi proyecto con IA responsable}

\begin{softbox}{milcGradeDark}{milcGradeSoft}{Saber ancestral}
\textbf{En la vereda La Plata de Cartago, cuando un aprendiz del taller de don Lucho terminaba su primer reloj reparado, lo presentaba a la comunidad.} No era trámite: era \emph{rito de paso}. El aprendiz mostraba el reloj funcionando, explicaba qué había encontrado averiado, qué herramientas usó, qué aprendió. Los vecinos hacían preguntas: \emph{``¿Por qué cambiaste esta pieza y no aquella? ¿Qué pasaría si fallara otra vez?''}. El aprendiz respondía. Si lo hacía bien, \textbf{la comunidad lo reconocía como artesano formado}. Esa transición pública era importante: \emph{no basta saber hacer; hay que mostrar el oficio}. Hoy es tu cosecha del periodo 3 y de todo el año de Tecnología. Hace 30 sesiones no sabías programar en Scratch, no entendías la nube, no sabías qué era un algoritmo, no habías hablado con una IA. \emph{Hoy entregas un proyecto que demuestra todo}. La comunidad académica (tu profe, tus compañeros) será testigo del oficio que ahora tienes.
\end{softbox}

\begin{softbox}{milcTurquesa}{milcGris}{Saber contemporáneo}
Esta es la \textbf{cosecha del periodo 3 y del año académico de grado 7°}. No es contenido nuevo: es \emph{integración}. Tu proyecto debe demostrar que sabes usar IA con criterio aplicando todo lo aprendido. \textbf{Las 5 dimensiones que debe tener:} \textbf{(1) Tema sólido:} algo informativo (no opinión personal), que TÚ escojas. Ejemplos: \emph{``Los volcanes del Eje Cafetero''}, \emph{``Historia del fútbol colombiano''}, \emph{``La cocina tradicional del Valle''}, \emph{``El cambio climático en Colombia''}. \textbf{(2) Voz tuya dominante:} no es que la IA escriba todo; tu voz, tus selecciones, tu organización dominan. \textbf{(3) IA como asistente declarado:} usaste IA y lo dices abiertamente, con plantilla de atribución. \textbf{(4) Anexo de transparencia:} muestras los prompts que usaste, las respuestas que recibiste, cómo verificaste, cómo adaptaste. \textbf{(5) Reflexión ética}: cómo cumpliste tus 5 acuerdos del uso responsable. \textbf{Es el proyecto más ambicioso del año --- y el más formativo}.
\end{softbox}

\begin{softbox}{milcMostaza}{milcGris}{Pregunta-puente}
\textit{Hace 30 sesiones (al iniciar grado 7°) \emph{no sabías nada de IA}. Hoy entregas un proyecto donde usas IA con criterio ético. \textbf{¿Cómo te sientes con ese crecimiento}?}
\end{softbox}

\begin{softbox}{milcVerde}{milcGris}{Saber y saber hacer}
Al terminar la clase: (1) podrás \textbf{integrar} todas las habilidades del periodo 3 en un proyecto; (2) sabrás \textbf{aplicar} prompting + ética + atribución en un trabajo real; (3) podrás \textbf{evaluar} tu propio uso de IA con rúbrica; (4) habrás \textbf{creado} tu proyecto + sustentado al profe.
\end{softbox}



\small
\begin{tabularx}{\textwidth}{>{\bfseries}p{3.0cm}Y}
\rowcolor{milcGradePrimary!12}Autor & Dr. Álvaro Cárdenas Orozco\\
DBA & Reconoce los fundamentos, usos y límites éticos de la Inteligencia Artificial (MEN, T\&I 7°)\\
Referentes & Saber ancestral del consejero · Modelos de lenguaje · Ética algorítmica y prompting\\
Producto final & Un \textbf{proyecto completo} que use IA con criterio: \textbf{(a) tema escogido} (informativo, no opinión); \textbf{(b) documento final} (Word o presentación) con tu voz dominante, IA como asistente declarado; \textbf{(c) anexo de transparencia}: 3-5 prompts CTRF/avanzados usados + capturas de las respuestas + reflexión sobre cómo verificaste/adaptaste; \textbf{(d) sustentación oral} de 5 minutos al profe; \textbf{(e) autoevaluación} con rúbrica de 10 criterios y \textbf{cierre del año}: ¿qué tipo de estudiante de IA soy ahora vs hace 10 sesiones?.\\
\end{tabularx}
\normalsize

\puente{Hoy haces 4 cosas: defines el tema, produces el proyecto con IA responsable, autoevalúas con rúbrica, sustentas.}

\milcestacion{milcGradeDark}{Ruta MILC: así vamos a aprender}
\begin{tabularx}{\textwidth}{>{\centering\arraybackslash}X>{\centering\arraybackslash}X>{\centering\arraybackslash}X>{\centering\arraybackslash}X}
\phasecard{1}{milcVerde}{milcGris}{Escucha\\[-1mm]\small Observo, escucho y pregunto.} &
\phasecard{2}{milcTurquesa}{milcGris}{Entender\\[-1mm]\small Sistematizo y ordeno.} &
\phasecard{3}{milcMagenta}{milcGris}{Hacer\\[-1mm]\small Construyo y aplico.} &
\phasecard{4}{milcVino}{milcGris}{Cierre\\[-1mm]\small Evalúo y reflexiono.}
\end{tabularx}


\puente{Empezamos definiendo el tema con las 4 técnicas del pensamiento computacional (de P2).}

\milcestacion{milcVerde}{Estación 1 de 3 · Escucha}

\begin{softbox}{milcVerde}{milcGris}{Escena para escuchar}
\textbf{Actividad 1 · ANALIZA --- Define el tema con las 4 técnicas} (12 min · individual). Aplica las 4 técnicas del pensamiento computacional (P2·S3) a definir tu proyecto. En el cuaderno responde: \emph{(1) Mi tema escogido (1 frase clara, informativo)}. \emph{(2) DESCOMPOSICIÓN: ¿en qué 4-5 sub-temas se divide?}. \emph{(3) PATRONES: ¿qué sub-temas son similares?}. \emph{(4) ABSTRACCIÓN: ¿qué detalles dejo afuera por ahora?}. \emph{(5) ALGORITMO: orden general (entrada-proceso-salida)?}.
\end{softbox}

{\sffamily\bfseries\fontsize{8}{10}\selectfont\textcolor{milcNegro!55}{MARCA CUANDO LO HAYAS HECHO}}
\milccheckrow{Definí tema concreto e informativo.}
\milccheckrow{Apliqué las 4 técnicas del pensamiento computacional.}
\milccheckrow{Tengo plan claro antes de empezar a producir.}

\begin{infoband}{milcVerde}


\textbf{Lo que va al cuaderno (Actividad 1 · ANALIZA).} Título: «Actividad 1 --- Plan de mi proyecto P3». \textbf{Formato:} 5 respuestas estructuradas. \textbf{Extensión:} 1/2 página. \textbf{Sabes que terminaste cuando} tu plan está claro antes de tocar la IA.
\end{infoband}

\puente{Después aprendes la rúbrica de 10 criterios + estructura del proyecto.}

\milcestacion{milcTurquesa}{Estación 2 de 3 · Entender}

\begin{softbox}{milcTurquesa}{milcGris}{Lo que vas a entender}
Ahora conoce la rúbrica con la que se evaluará tu proyecto + estructura concreta del entregable.
\end{softbox}

\milcpaso{1}{milcTurquesa}{\textbf{Los 10 criterios de la rúbrica (cada uno 1-5 puntos, total 50).} \textbf{(1) Plan inicial} con las 4 técnicas aplicadas. \textbf{(2) Tema sólido} (informativo, acotado). \textbf{(3) Voz propia dominante} (tu organización, tus selecciones, tu cierre). \textbf{(4) IA usada con criterio} (asistente, no reemplazo). \textbf{(5) Atribución correcta} (plantilla aplicada). \textbf{(6) Anexo de transparencia} (prompts + respuestas + verificación). \textbf{(7) 3+ fuentes citadas} (no solo IA; también web/libros). \textbf{(8) Cumplimiento de los 5 acuerdos} (verificar, no datos personales, citar, pensar, voz propia). \textbf{(9) Sustentación clara} de 5 minutos. \textbf{(10) Reflexión final} sobre crecimiento del año.}
\milcpaso{2}{milcTurquesa}{\textbf{Estructura del entregable.} (a) \textbf{Documento principal} en Word (o presentación en PowerPoint): 2-3 páginas, formato profesional. (b) \textbf{Anexo de transparencia}: hoja aparte con 3-5 prompts CTRF/avanzados usados + capturas de respuestas + 1 línea de cómo verificaste o adaptaste cada uno. (c) \textbf{Reflexión final}: 1 párrafo sobre tu crecimiento. (d) \textbf{Sustentación oral}: 5 minutos donde muestras el proyecto, explicas decisiones, atribución, ética. El profe puede preguntar.}
\milcpaso{3}{milcTurquesa}{\textbf{3 cosas que diferencian un proyecto básico de uno excelente.} (1) \emph{Originalidad:} tema único + voz propia + selección personal del ángulo. (2) \emph{Coherencia ética:} no solo se usaron las herramientas, se aplicó la ética con criterio. (3) \emph{Honestidad transparente:} el anexo muestra abiertamente cómo se usó la IA. Un profesor experimentado lo notará y lo valorará.}
\milcpaso{4}{milcTurquesa}{\textbf{4 errores típicos en la cosecha P3.} (1) \emph{Querer hacer el proyecto sin IA para evitar el debate ético} --- pero la S10 te pide usarla con criterio. (2) \emph{Pegar la respuesta de ChatGPT sin verificar} --- las alucinaciones se cuelan. (3) \emph{Olvidar la atribución} --- viola el acuerdo 3. (4) \emph{No reflejar tu voz} --- todo suena a IA, sin tu personalidad. La cosecha exige tu voz dominante.}

\begin{drawbox}{milcTurquesa}{10 criterios + estructura del entregable}
\textbf{10 criterios:} plan · tema sólido · voz propia · IA con criterio · atribución · anexo transparencia · fuentes · 5 acuerdos · sustentación · reflexión. \\[1mm]
\textbf{Entregable:} documento principal + anexo transparencia + reflexión + sustentación.
\end{drawbox}

\begin{softbox}{milcMostaza}{milcGris}{Errores comunes y su corrección}
\textbf{Mitos sobre la cosecha P3.} (1) \emph{``Si no uso IA, evito problemas''} --- Falso. La cosecha pide usarla con criterio, no evadirla. (2) \emph{``Si la voz de la IA es buena, da igual cuál es la mía''} --- Falso. El profe quiere tu voz, no la de ChatGPT. (3) \emph{``El anexo es opcional''} --- No. Es 10\% de la rúbrica. (4) \emph{``La sustentación es lo de menos''} --- Falso. Es 10\% directo. (5) \emph{``Si entrego ya, no necesito reflexión''} --- Falso. La reflexión cierra el año académico.
\end{softbox}

\begin{infoband}{milcTurquesa}


\textbf{Lo que va al cuaderno (Actividad 2 · EXPLICA).} Título: «Actividad 2 --- Rúbrica + estructura del entregable». \textbf{Formato:} lista de 10 criterios + esquema de estructura del entregable. \textbf{Extensión:} 15 líneas. \textbf{Sabes que terminaste cuando} puedes autoevaluar cualquier proyecto con la rúbrica.
\end{infoband}

\puente{Con todo claro, produces el proyecto + anexo de transparencia + sustentación.}

\milcestacion{milcMagenta}{Estación 3 de 3 · Hacer}

\begin{softbox}{milcMagenta}{milcGris}{Cómo vas a construir tu producto}
\textbf{Actividad 3 · CREA --- Proyecto completo con IA responsable} (1.5 sesiones de 90 min + trabajo asíncrono durante la semana). Vas a producir el proyecto cosecha. Es la culminación del año académico de grado 7°.
\end{softbox}

\milcpaso{1}{milcMagenta}{\textbf{Paso 1 --- Documento principal (40-50 min).} En Word (o Docs) crea un documento con tu tema. Estructura: \emph{introducción (1 párrafo), desarrollo (4-5 párrafos con tus subtemas), conclusión (1 párrafo)}. Total 2-3 páginas con formato profesional (negrita en ideas clave, sangría, fuentes citadas al pie). \textbf{Importante: tu voz domina.} Si usas IA para ayudar, adáptala a tu manera de escribir, tus ejemplos, tu perspectiva colombiana. Guarda como \texttt{2026-cosecha-p3-[tu-tema]-[tu-nombre]}.}
\milcpaso{2}{milcMagenta}{\textbf{Paso 2 --- Anexo de transparencia (15-20 min).} Crea un documento aparte (puede ser Word, puede ser parte del mismo): \textbf{``Anexo --- Cómo usé IA''}. Para cada uso de IA en el proyecto, escribe: \emph{(a) Prompt usado completo (CTRF + técnicas avanzadas)}, \emph{(b) Resumen de la respuesta recibida (3-4 líneas)}, \emph{(c) Cómo verifiqué o adapté esa respuesta para incluirla en mi proyecto}. Mínimo 3 ejemplos. Esto demuestra al profe que usaste IA con criterio, no la copiaste.}
\milcpaso{3}{milcMagenta}{\textbf{Paso 3 --- Atribución dentro del documento.} En la última página del documento principal, sección \textbf{``Fuentes y asistencia''}: lista las fuentes humanas (web, libros, entrevistas) Y la asistencia de IA con la plantilla: \emph{``En este proyecto usé [LLM] para [propósito específico]. Las ideas, organización, voz y conclusiones son mías.''}.}
\milcpaso{4}{milcMagenta}{\textbf{Paso 4 --- Reflexión final.} Sección aparte (puede ser en el cuaderno): \textbf{``Mi crecimiento del año''}. Escribe en 5-7 líneas: \emph{(a) ¿Qué sabía al inicio del año en grado 7°?}, \emph{(b) ¿Qué sé al final del año?}, \emph{(c) ¿Cuál fue el aprendizaje que más me sorprendió?}, \emph{(d) ¿Cómo me sentí al usar IA por primera vez con criterio ético?}, \emph{(e) ¿Qué le recomendaría a mi yo de hace 30 sesiones?}.}
\milcpaso{5}{milcMagenta}{\textbf{Paso 5 --- Autoevaluación con rúbrica.} Antes de sustentar, abre la rúbrica de 10 criterios. Marca cada uno con tu nota (1-5). Suma. Si algún criterio quedó bajo, ajusta si tienes tiempo. \textbf{Honesta autoevaluación es señal de madurez adulta}.}
\milcpaso{6}{milcMagenta}{\textbf{Paso 6 --- Sustentación de 5 minutos al profe (en clase).} Estructura sugerida: \emph{(a) 30s: presento el tema, por qué lo elegí}. \emph{(b) 2 min: muestro el documento principal, explico organización}. \emph{(c) 1 min: muestro anexo de transparencia, explico cómo usé IA con criterio}. \emph{(d) 30s: leo mi reflexión final del año}. \emph{(e) 1 min: preguntas del profe}. La sustentación cierra el oficio.}

\begin{drawbox}{milcMagenta}{Antes de sustentar}
\checkbox Documento principal con tema sólido y voz propia dominante. \\[1mm]
\checkbox Anexo de transparencia con 3+ prompts y verificaciones. \\[1mm]
\checkbox Atribución de IA correctamente aplicada. \\[1mm]
\checkbox 3+ fuentes humanas citadas además de IA. \\[1mm]
\checkbox Reflexión final del año del crecimiento personal. \\[1mm]
\checkbox Autoevaluación honesta con rúbrica de 10 criterios.
\end{drawbox}

\begin{softbox}{milcMagenta}{milcGris}{Plantilla de trabajo}
\textbf{Estructura.} \textit{Paso 1:} documento principal 2-3 páginas con voz propia. \textit{Paso 2:} anexo de transparencia con 3+ prompts. \textit{Paso 3:} atribución de IA. \textit{Paso 4:} reflexión final del año. \textit{Paso 5:} autoevaluación con rúbrica. \textit{Paso 6:} sustentación oral 5 min.
\end{softbox}

\begin{infoband}{milcMagenta}


\textbf{Lo que va al cuaderno (Actividad 3 · CREA).} Título: «Actividad 3 --- Cosecha P3 + cierre del año». \textbf{Formato:} bocetos + anexo de transparencia + reflexión final firmada. \textbf{Extensión:} 1.5 a 2 páginas. \textbf{Sabes que terminaste cuando} sustentaste el proyecto al profe.
\end{infoband}

\puente{El producto es proyecto + anexo + reflexión + sustentación al profe.}

\milcestacion{milcMostaza}{Producto de la sesión}
\textbf{Proyecto con IA responsable + anexo de transparencia + sustentación · cierre del año académico}.
\begin{softbox}{milcMostaza}{milcGris}{Criterios de calidad}
(1) El documento principal tiene tema sólido, voz propia y formato profesional. (2) El anexo de transparencia documenta 3+ usos de IA con verificaciones. (3) La atribución de IA está aplicada correctamente. (4) La reflexión final del año académico es honesta y profunda. (5) La sustentación oral fue clara y demostrativa.
\end{softbox}

\puente{Antes de sustentar, autoevalúa con la rúbrica de 10 criterios.}

\milcestacion{milcVino}{Evaluación liberadora · cinco dimensiones}

\begin{softbox}{milcVino}{milcGris}{Sentido de la evaluación}
Evaluar no es solo revisar una nota. Es reconocer qué aprendí, qué cambió en mí, qué emociones manejé, cómo me relaciono con la ciudad y la comunidad, y qué le dejaré a quien viene detrás.
\end{softbox}

\textbf{Mi reflexión liberadora en cinco dimensiones.} Completa cada frase iniciada:

\small
\begin{tabularx}{\textwidth}{>{\bfseries\RaggedRight\arraybackslash}p{3.4cm}Y>{\RaggedRight\arraybackslash}p{4.6cm}}
\rowcolor{milcVino}\color{white}Dimensión & \color{white}\bfseries Mi respuesta (frame starter) & \color{white}\bfseries Algo que descubrí\\
1. Personal & Lo que más cambió en mí fue \linead{6.5cm} & \linead{4.2cm}\\
2. Emocional & La emoción más fuerte fue \linead{2.5cm} y la regulé \linead{3cm} & \linead{4.2cm}\\
3. Ciudadana & En democracia esto importa porque \linead{6cm} & \linead{4.2cm}\\
4. Local (Cartago) & En mi barrio sería útil para \linead{6.5cm} & \linead{4.2cm}\\
5. Intergeneracional & A mi abuela/abuelo le diría \linead{6.5cm} & \linead{4.2cm}\\
\end{tabularx}
\normalsize

\begin{infoband}{milcVino}
\textbf{Para la próxima sustentación.} Vuelve a esta tabla en seis meses. Verás que las respuestas cambian — no porque lo aprendido cambie, sino porque tú cambias. La evaluación liberadora es una conversación contigo a través del tiempo.
\end{infoband}


\puente{Cerramos con 3 voces sobre el cierre del año y el camino hacia 8° grado.}

\milcestacion{milcGradeDark}{Triángulo de pensamiento · cierre filosófico}

\begin{softbox}{milcVino}{milcGris}{Dussel · lente del nosotros}
\textit{``Aprender a usar IA con criterio ético en grado 7° es entrar a la adultez digital antes de tiempo. Es ventaja real, no solo académica.''}

\textbf{La mayoría de adultos hoy NO saben usar IA con criterio ético}. Tu generación entra al mundo digital con esa habilidad como base. \emph{Eso te coloca en lugar privilegiado para la vida adulta}: en universidad, en trabajo, en proyectos sociales. No es exageración: \emph{conocer y dominar las herramientas que tu generación va a usar profesionalmente es invertir en tu yo de 25 años}.

\textbf{¿Reconozco el valor de haber aprendido esto en momento clave?}
\end{softbox}
\linead{\linewidth}\\[1mm]
\linead{\linewidth}

\begin{softbox}{milcMostaza}{milcGris}{Estoicismo · Marco Aurelio · cuidado interior}
\textit{``Cerrar un año con oficio cumplido es base para el siguiente. Cierra bien grado 7°; abrirás bien grado 8°.''}

\textbf{Hoy cierras grado 7° en Tecnología}. Si lo cierras bien --- proyecto entregado con criterio, sustentación clara, reflexión honesta --- \emph{abres grado 8° con base sólida}. Si lo cierras a medias o sin oficio, en 8° empiezas desde atrás. \emph{El estoico sabe que cada cierre prepara el siguiente comienzo}. Tu cosecha hoy es promesa de un buen 8° grado en 2027.

\textbf{¿Estoy cerrando este año con el oficio que quiero llevar a 8°?}
\end{softbox}
\linead{\linewidth}\\[1mm]
\linead{\linewidth}

\begin{softbox}{milcTurquesa}{milcGris}{Floridi · lente de la infoesfera}
\textit{``El primer proyecto con IA responsable es entrada formal a la ciudadanía digital adulta. No hay vuelta atrás: ya eres ciudadano.''}

\textbf{Hoy cambias de identidad}: antes eras \emph{usuario novato de tecnología}; ahora eres \emph{ciudadano digital con criterio ético}. Ese cambio se queda. \emph{En 30 años más vas a recordar} este año de grado 7° como el momento donde aprendiste a usar IA con criterio. Tu yo adulto te va a agradecer haber prestado atención. Sigue la línea: cada año en colegio, secundaria, universidad, trabajo, vas a profundizar lo que empezaste aquí.

\textbf{¿Reconozco que hoy soy un ciudadano digital con criterio, no un usuario pasivo más?}
\end{softbox}
\linead{\linewidth}\\[1mm]
\linead{\linewidth}

\begin{infoband}{milcNegro}
\textbf{Nota para el docente.} Las tres formulaciones del triángulo son la síntesis del autor de la guía, no citas textuales de los pensadores. Si vas a citarlos en clase, remite a la obra original.
\end{infoband}

\milcestacion{milcVino}{Mi autoevaluación · marca una casilla por fila}

{\small
\setlength{\extrarowheight}{2.2pt}
\begin{tabularx}{\textwidth}{>{\RaggedRight\arraybackslash\bfseries}p{3.0cm}YYY}
\rowcolor{milcVino}\color{white}\bfseries Criterio & \color{white}\bfseries Lo logré & \color{white}\bfseries En proceso & \color{white}\bfseries Necesito apoyo\\[.6mm]
Comprensión del tema & \milcopcion{Explico las ideas con mis palabras.} & \milcopcion{Reconozco las ideas, falta precisión.} & \milcopcion{Necesito repasar lo básico.}\\[.8mm]
\rowcolor{milcGris}Pensamiento computacional & \milcopcion{Usé los pilares al armar mi producto.} & \milcopcion{Usé algunos pilares.} & \milcopcion{Necesito releer la estación 2.}\\[.8mm]
Aplicación y producto & \milcopcion{Mi producto cumple los criterios.} & \milcopcion{Está armado, con ajustes pendientes.} & \milcopcion{Aún está en construcción.}\\[.8mm]
\rowcolor{milcGris}Reflexión 5 dimensiones & \milcopcion{Respondí cada una con honestidad.} & \milcopcion{Respondí algunas con detalle.} & \milcopcion{Mis respuestas son muy generales.}\\[.8mm]
Triángulo de pensamiento & \milcopcion{Conecté las 3 voces con mi trabajo.} & \milcopcion{Conecté al menos una.} & \milcopcion{No las relaciono con mi proceso.}\\[.8mm]
\end{tabularx}}

\vspace{3mm}
\textbf{Hoy entendí que:}\\[1mm]
\linead{\linewidth}\\[1mm]
\linead{\linewidth}

\textbf{Mi compromiso para la próxima sesión es:}\\[1mm]
\linead{\linewidth}\\[1mm]
\linead{\linewidth}

\begin{softbox}{milcMostaza}{milcGris}{Cita de despedida · Marco Aurelio}
\textit{``El obstáculo en el camino se vuelve el camino. Lo que estorba al camino se vuelve camino.''}\\[1mm]
Cada error de hoy, bien observado, se convierte en pieza del camino que pisarás mañana.
\end{softbox}

\small
\begin{tabularx}{\textwidth}{>{\bfseries}p{3.6cm}Y}
\rowcolor{milcGradeDark!12}Nombre completo: & \linead{10cm}\\
Fecha: & \linead{4cm} \quad Curso/grupo: \linead{4cm}\\
Mi firma: & \linead{9cm}\\
\end{tabularx}
\normalsize

\begin{infoband}{milcGradeDark}
\textbf{Para la siguiente sesión.} Guarda tu proyecto en lugar seguro. La S11 es el \textbf{cierre formal de grado 7°}: autoevaluación con rúbrica visible del año, examen integrador, mirada a grado 8° que se viene en 2027 (Excel, lógica de control, multimedia). Hoy cerraste el año académico con tu cosecha. La próxima clase consolidas formalmente todo lo que aprendiste.
\end{infoband}

\end{document}
